%==============================================================================
%  27-glossary.tex — Glossary
%==============================================================================
\strategyPart{27 \textbar\ Glossary}{Glossary of Key Terms}

A compact glossary for consistent use of terms across this survey and any
documents derived from it [ISO/IEC 22989; NIST AI RMF].

\begin{description}
  \item[Agent (AI agent)] A system that receives an objective, plans steps,
        uses tools and may execute actions with limited human intervention.
  \item[AI-native organization] An organization whose products, operations
        and decision-making are designed around AI from the start.
  \item[AI-orchestrated organization] An organization in which agents
        coordinate information, tasks, decisions and workflows across teams
        and systems.
  \item[Assistant (AI assistant)] A system that supports a person by
        retrieving, summarizing, recommending or drafting, without acting
        independently.
  \item[Autonomy level] A graded scale (0--6) describing how much of a task
        an AI system performs and how much human involvement remains.
  \item[Copilot] An AI that works inside a business process and suggests
        actions while the employee remains responsible for the final action.
  \item[Data capital] The stock of trusted, governed, accessible data an
        organization can use.
  \item[Foundation model] A large pretrained model (e.g., an LLM) that can be
        adapted to many tasks.
  \item[Guardrails] Technical and organizational controls that constrain AI
        behavior (permissions, filters, approvals).
  \item[Human capital] The stock of AI-literate talent, leadership and
        workforce skills.
  \item[Human-in-the-loop] A design in which a human approves, reviews or
        monitors AI actions at defined checkpoints.
  \item[LLMOps] Operational practices for running large language models in
        production (routing, cost, evaluation, monitoring).
  \item[MLOps] Operational practices for the model lifecycle (registry,
        deployment, monitoring, retraining).
  \item[Multi-agent system] Several specialized agents that cooperate (e.g.,
        planning, retrieval, analysis, validation, reporting, governance).
  \item[Organizational capital] The stock of redesigned workflows,
        incentives, decision rights and collaboration structures.
  \item[Proof of concept (PoC)] A demonstration that an AI use case can work,
        before production scaling.
  \item[RAG] Retrieval-Augmented Generation: grounding model outputs in
        retrieved, trusted sources with citations.
  \item[Shadow AI] Use of unapproved AI tools or data outside governance.
  \item[Workflow (AI workflow)] A predefined, engineer-controlled sequence of
        AI steps.
\end{description}
